\section{Fourth-order value asymptotics}\label{sec:general}
At each simple zero \(z_k\) of \(q_*\), write
\[
 Q_k=(q_j(z_k))_{j<m},\quad r_k=q_*'(z_k),\quad
 t_k=q_*''(z_k),\quad u_k=q_*'''(z_k),\quad
 \gamma_k=|r_k|,\quad \sigma_k=\sgn r_k.
\]
Define the following quantities whenever their series converge:
\begin{align}
 \Gamma&=\sum_k\gamma_k,&
 G&=2\sum_k\frac{Q_kQ_k^{\mathsf T}}{\gamma_k},\label{eq:GB}\\
 B_j&=\frac13\sum_k\sigma_k
       \left(\frac{q_j(z_k)t_k}{r_k}-q_j'(z_k)\right),&
 v&=G^{-1}B,\nonumber\\
 R&=\sum_k\left(\frac{t_k^2}{36\gamma_k}
                    -\frac{\sigma_ku_k}{60}\right),&
 P&=\frac12B^{\mathsf T}G^{-1}B,\qquad \Xi=R-P.\label{eq:RP}
\end{align}
Here \(G\) is an \(m\)-by-\(m\) matrix and \(B,v\in\R^m\).

\subsection{Countably many switches}
We impose the following hypotheses on \(X=[0,\infty)\).
\begin{enumerate}[label=(H\arabic*),leftmargin=*]
\item The functions \(q_j\) are in \(L^1(X)\) and are locally \(C^3\).
The sign certificate \eqref{eq:signcertificate} holds, \(q_*(0)\ne0\), and
the positive zeros \(z_k\) tend to infinity with positive minimum separation.
\item There are disjoint neighborhoods \(N_k=[z_k-\rho,z_k+\rho]\subset(0,\infty)\),
with a common \(\rho>0\), and a bounded open neighborhood \(U\) of \(b^*\).
For every \(b\in U\), \(q_b\) has exactly one zero \(z_k(b)\) in each \(N_k\)
and no other zeros. The orientations remain \(\sigma_k\).
The functions \(z_k(b)\) are \(C^1\), have a common Lipschitz constant, and
\(|z_k(b)-z_k|<\rho/4\).
\item There are \(L_k>0\) such that
\(|q_b(x)|\geq L_k|x-z_k(b)|\) on \(N_k\) for all \(b\in U\).
Set
\[
 E_k=\sum_{j=0}^m\sum_{\ell=0}^3\sup_{N_k}|q_j^{(\ell)}|,
 \quad C_U=1+\sup_{b\in U}\sum_{j<m}|b_j|,
 \quad K_k=\frac{C_UE_k}{6L_k}.
\]
The weighted condition is
\begin{equation}\label{eq:weighted}
 \sum_kE_k(1+K_k)^2<\infty.
\end{equation}
The sequence \(K_k\) need not be bounded.
\item Some \(m\) vectors among the \(Q_k\) are linearly independent.
\end{enumerate}

\begin{theorem}\label{thm:infinite}
Under {\rm(H1)--(H4)}, all the series in \eqref{eq:GB}--\eqref{eq:RP}
converge absolutely, \(G\) is positive definite, and \eqref{eq:problem}
has an optimizer for all sufficiently large \(M\). Moreover,
\begin{align}
 \delta(M)&=\delta_0+C_2M^{-2}+C_4M^{-4}+o(M^{-4}),\label{eq:expansion}\\
 C_2&=\frac{\delta_0^3\Gamma}{3D},&
 C_4&=\delta_0^5\left(\frac{\Gamma^2}{3D^2}-\frac{R}{D}
                          +\frac{P}{D}\right).\label{eq:coefficients}
\end{align}
The term \(P\) is nonnegative and is unchanged by any invertible linear
change of basis of the lower moment functions \(q_0,\ldots,q_{m-1}\).
\end{theorem}

We give the proof, including the tail and exact-feasibility steps.
For an average over \([-1,1]\), write
\(\mathbb E h(y)=\frac12\int_{-1}^1h(y)\dd y\).

\begin{lemma}[Balanced ramps]\label{lem:ramps}
For all \(b\in U\) sufficiently close to \(b^*\) and all sufficiently small
\(a>0\), there is a unique center \(c_k\in(z_k(b)-a,z_k(b)+a)\) satisfying
\begin{equation}\label{eq:balance}
 \mathbb E q_b(c_k+ay)=0.
\end{equation}
Its displacement satisfies
\[
 |c_k-z_k(b)|\leq\min\{a,K_ka^2\}.
\]
The function equal to \(\sgn q_b\) outside the ramp intervals and equal
to \(\sigma_k(x-c_k)/a\) on \([c_k-a,c_k+a]\) maximizes \(H_b(a)\).
\end{lemma}
\begin{proof}
At \(c=z_k(b)-a\) and \(c=z_k(b)+a\), the averaged function has opposite
signs. Between them its derivative is
\([q_b(c+a)-q_b(c-a)]/(2a)\), of strict sign \(\sigma_k\).
Taylor's theorem gives
\(|q_b(c)-\mathbb E q_b(c+ay)|\leq C_UE_ka^2/6\).
The crossing bound then proves the displacement estimate.

On a balanced interval \([\ell,r]\), put
\(J(x)=\int_\ell^xq_b(t)\dd t\).
There is just one change of sign of \(q_b\), and
\(J(\ell)=J(r)=0\), \(\sigma_kJ(x)<0\) inside.
For any competitor \(s\), integration by parts gives
\[
 \int_\ell^rq_b(s-s_a)
   =-\int_\ell^rJ(x)(s'(x)-\sigma_k/a)\dd x\leq0.
\]
Outside these intervals, \(s_a=\sgn q_b\) maximizes the integrand
pointwise. The intervals are disjoint for small \(a\).
The infinite sum is justified by the integrable bound \(2|q_b|\).
\end{proof}

\begin{proof}[Proof of Theorem~\ref{thm:infinite}]
\emph{Convergence and coefficients.}
Differentiating the zero equation gives
\(\partial_{b_j}z_k=q_j(z_k)/r_k\).
The common root Lipschitz bound in (H2) therefore bounds \(Q_k/r_k\).
Also \(|t_k/r_k|\leq6K_k\).
These estimates and \eqref{eq:weighted} imply absolute convergence of
\(\Gamma,G,B,R\). Condition (H4) makes \(G\) positive definite.

Choose the trial multiplier \(b(a)=b^*+va^2\) and its balanced ramp
profile \(s_a\). This trial multiplier is not asserted to solve the exact
dual optimization problem. Define
\begin{equation}\label{eq:kappa}
 \kappa_k=\frac{Q_k^{\mathsf T}v-t_k/6}{r_k}.
\end{equation}
Implicit differentiation of the zero equation and
\eqref{eq:balance}, for each fixed \(k\), give
\[
 \Delta_k:=c_k-z_k=\kappa_ka^2+o(a^2).
\]
Uniformly in \(k\), for a fixed constant \(C\),
\begin{equation}\label{eq:displacements}
 |\Delta_k|\leq a+C\norm{v}a^2\leq2a,\qquad
 |\Delta_k|/a^2\leq K_k+C\norm{v}.
\end{equation}
The first estimate and the second serve different purposes; we do not
replace them by a false uniform bound on \(\kappa_k\).

\emph{Summed moment expansions.}
For \(H_j(t)=\int_t^\infty q_j(x)\dd x\), the moment change at a switch is
\[
 2\sigma_k\{\mathbb E H_j(z_k+\Delta_k+ay)-H_j(z_k)\}.
\]
Only these differences are summed; separate sums of the \(H_j(z_k)\)
are unnecessary. If \(h=\Delta_k+ay\), then
\[
 \mathbb E h=\Delta_k,\quad
 \mathbb E h^2=\Delta_k^2+a^2/3,\quad
 \mathbb E h^3=\Delta_k^3+\Delta_ka^2,\quad
 \mathbb E h^4=\Delta_k^4+2\Delta_k^2a^2+a^4/5 .
\]
Taylor expansion and \eqref{eq:displacements} now yield
\begin{align}
 \int q_js_a&=a^2(B-Gv)_j+o(a^2)=o(a^2)\quad(j<m),\label{eq:defect}\\
 \int q_*s_a&=D-\frac{\Gamma a^2}{3}
 -a^4\sum_k\sigma_k
 \left(r_k\kappa_k^2+\frac{t_k\kappa_k}{3}+\frac{u_k}{60}\right)
 +o(a^4).\label{eq:targetexp}
\end{align}
Here the normalized lower moment terms and the target terms are dominated
by a constant times \(E_k(1+K_k)^2\).
For example, \(r_k\Delta_k^2/a^4\) uses the second bound in
\eqref{eq:displacements}; terms of degree at least four use
\(|\Delta_k|\leq2a\) as well. The Taylor remainder divided by \(a^4\)
tends to zero for each fixed \(k\), by continuity of \(q_*'''\), and has
the same summable majorant. Dominated convergence proves both displayed
expansions.

Completing the square with \eqref{eq:kappa} gives
\[
 -\sum_k\sigma_k
 (r_k\kappa_k^2+t_k\kappa_k/3+u_k/60)
   =R-\tfrac12v^{\mathsf T}Gv=\Xi .
\]
Since \(q_{b(a)}=q_*-a^2\sum v_jq_j\),
\eqref{eq:defect} also gives the support expansion
\begin{equation}\label{eq:supportexp}
 H_{b(a)}(a)=D-\Gamma a^2/3+\Xi a^4+o(a^4).
\end{equation}
For a general trial vector \(v\), its fourth coefficient would be
\(R+\tfrac12v^{\mathsf T}Gv-B^{\mathsf T}v\);
its minimum occurs at \(G^{-1}B\).

\emph{Exact moment repair.}
Select the \(m\) independent switches from (H4), and move only their
centers by a vector \(e\). At \(a=0,e=0\), the derivative of the lower
moment map has columns \(-2\sigma_kQ_k\), so it is invertible.
For small \(a\) this derivative remains uniformly close to that invertible
matrix. A contraction with that matrix as preconditioner, on a ball of
radius a fixed multiple of the defect in \eqref{eq:defect}, gives a
continuous solution \(e(a)=o(a^2)\).
Indeed its forcing is \(o(a^2)\), and its derivative norm can be made
less than \(1/2\), so that this ball maps into itself.
Ramp intervals remain disjoint and the amplitude and slope bounds are
unchanged. The derivative of the \(q_*\) moment with respect to each
selected center is \(O(a^2)\), while its second derivative is bounded.
Consequently the repair changes that moment by
\(O(a^2|e|+|e|^2)=o(a^4)\).
The repaired profile \(\widetilde s_a\) therefore has exact lower
moments and
\begin{equation}\label{eq:primalexp}
 S(a):=\int q_m\widetilde s_a
      =D-\Gamma a^2/3+\Xi a^4+o(a^4).
\end{equation}

Put \(A(a)=f/S(a)\) and \(M(a)=A(a)/a\).
The feasible function \(A(a)\widetilde s_a\) has amplitude \(A(a)\)
and slope at most \(M(a)\). Both functions are continuous for small
positive \(a\), \(A(a)\to\delta_0\), and \(M(a)\to\infty\).
The image of this interval under \(M(a)\) is an interval unbounded
above, hence contains all sufficiently large budgets. No monotonicity
of this parametrization is needed.

\emph{Attainment and matching lower bound.}
For a fixed sufficiently large \(M\), a minimizing sequence is uniformly
bounded and equicontinuous. A diagonal Arzel\`a--Ascoli argument gives a
locally uniform limit. Dominated convergence against each \(q_j\in L^1\)
preserves the moments, and the norm and Lipschitz bounds pass to the
limit. Thus the optimum is attained.

For any feasible amplitude \(A\), use
\(a=A/M\), \(s=g/A\), and the trial multiplier at that \(a\). Then
\[
 f=\int q_{b(a)}g\leq A H_{b(a)}(a).
\]
The preceding feasible upper bound and \(\delta(M)\geq\delta_0\)
show that the minimizing amplitude tends to \(\delta_0\).
The support and primal expansions therefore have matching scalar
inversions. Explicitly, with
\(t=\Gamma/(3D)\), \(e=\Xi/D\), the equation to this order is
\[
 A\{1-t(A/M)^2+e(A/M)^4+o(M^{-4})\}=\delta_0 .
\]
Its increasing polynomial part has derivative bounded away from zero
near \(\delta_0\). The remainder is uniform for \(A\) in a compact
neighborhood of \(\delta_0\). Substitution yields
\[
 A=\delta_0+t\delta_0^3M^{-2}
                   +(3t^2-e)\delta_0^5M^{-4}+o(M^{-4}),
\]
which proves \eqref{eq:coefficients}. In particular the factor \(3\)
accounts for the dependence \(a=A/M\).
Finally, a lower-moment basis change \(Q\mapsto LQ\) gives
\(G\mapsto LGL^{\mathsf T}\), \(B\mapsto LB\), and leaves \(P\) unchanged.
\end{proof}

\subsection{Finite switch sets and the role of the constraints}
\begin{proposition}\label{prop:finite}
On a compact interval, suppose the \(q_j\) are \(C^2\), that
\eqref{eq:signcertificate} holds, that all zeros of \(q_*\) are finitely
many simple interior zeros, with no endpoint zeros, and that the \(Q_k\)
span \(\R^m\). For all sufficiently large \(M\), the optimum in
\eqref{eq:problem} is unique and consists of constant plateaus joined by
linear ramps. Its value satisfies
\[
 \delta(M)=\delta_0+\frac{\delta_0^3\Gamma}{3D}M^{-2}+o(M^{-3}).
\]
If \(q_*\) is \(C^3\) near the zeros, the full expansion
\eqref{eq:expansion}--\eqref{eq:coefficients} holds with finite sums.
\end{proposition}
\begin{proof}
Root separation and sign persistence are automatic on a compact interval
after shrinking the neighborhood of \(b^*\).
The balance equations define centers \(c_k(a,b)\).
At \(a=0\) the lower-moment map of their ramp profile has \(b\)-derivative
\(-G\). Expressing its local contributions as averages of primitives
shows that this map is \(C^2\), including at \(a=0\).
The implicit function theorem supplies an even \(C^2\) multiplier \(b(a)\)
which makes all lower moments exactly zero.
Thus \(b(a)=b^*+va^2+o(a^2)\), with \(v\) as above.

Lemma~\ref{lem:ramps} proves exact support optimality.
Equality forces \(s=\sgn q_b\) on every plateau and
\(s'=\sigma_k/a\) almost everywhere in every ramp, because
\(\sigma_kJ<0\) there. Continuity fixes all endpoint values.
The support maximizer, and hence the constrained optimizer at the
resulting budget, is unique.
The target satisfies \(S(a)=D-\Gamma a^2/3+o(a^3)\) under \(C^2\)
regularity; its derivative is \(O(a)\).
Hence \(M(a)=f/(aS(a))\) has negative derivative for small \(a>0\)
and covers all sufficiently large budgets. The scalar inversion gives
the stated second-order result. If \(q_*\) is \(C^3\) near the switches,
the finite version of \eqref{eq:targetexp} gives the fourth coefficient.
\end{proof}

\begin{remark}
The sign of \(C_4\) is not fixed by positivity of an underlying weight.
The nonnegative quantity \(P\) has a more specific interpretation.
For the fixed residual \(q_*\), the unconstrained ramp support has
fourth coefficient \(R\); imposing the lower moments gives \(\Xi=R-P\).
Thus the loss at this order is \(Pa^4\).
This does not compare the original \(q_m\) problem with a different
problem in which all lower constraints have simply been deleted.
\end{remark}
